\subsection{The all-arrows theorem and the distinction between its models}

\begin{result}{E.1}{Known}{The verified all-arrows comparison}
Let $X$ be a small category and let $J$ be the integer fence with arrows $2i\to2i\pm1$. Write $\mathcal P(X)$ for the category of functors $J\to X$ constant with identity arrows outside a finite interval, and $p=(\sigma,\beta):\mathcal P(X)\to X\times X$ for the eventual endpoints. Then the constant-path section $X\to\mathcal P(X)$ is a weak equivalence, the one-endpoint fibres are weakly contractible, and the exact two-endpoint fibres are homotopy fibres. The canonical comparison span between these fibres and the all-type fibres $\Ch(X;a,b)$ consists of weak equivalences and respects composition in $\Hot$. Thus
\[
 B\Ch(X;a,b)\simeq B\mathcal P(X;a,b)\simeq P_{a,b}(BX).
\]
The actual even translation on $\mathcal P(X;a,b)$ is the identity in $\Hot$. The actual one-sided constant-tail maps $q_L,q_R$ are weak equivalences, including for any wide relative class on the corresponding relative fibres.
\end{result}
\begin{proof}[Verification and precise use of the supplied proofs]
We checked CATHOT E.20--E.25 and the companion's blocks at VII ms.~166--169, including the variance of the lifting categories. In accordance with the task, we cite those checked proofs rather than reproduce the all-arrows argument. The essential hypotheses are positive-length supports when invoking a finite-stage Grothendieck fibration, both endpoint identities on fibre morphisms, and filtered unions on objects and arrows. The all-type endpoint functor has only the smoothness used in the repaired argument; its properness is false, as E.3 records.

The shorter $q_L$ argument in the supplied second review is valid. Let $\mathcal P_{\leq2k}$ be the full subcategory of paths constant on $[2k,\infty)$. Its union is $\mathcal P$, and after translation each successive inclusion is the collar-shift map $S$ of CATHOT E.23. The explicit stationary-on-finite-support collar homotopy shows $BS\simeq1$. A sequential union of these weak equivalences is a weak equivalence because the nerve and homotopy groups after $\operatorname{Ex}^{\infty}$ commute with filtered colimits. This applies on the exact fibres and preserves the relative backward-arrow condition. It does not give a bounded zigzag of natural transformations; the counterexample in Part C shows why that stronger conclusion is false.
\end{proof}

\begin{result}{E.2}{Proved}{A strict categorical multiplication on the loop model}
For every $x\in X$, the all-type category $\Omega_{\mathrm{cat}}(X,x)=\Ch(X;x,x)$ is a strict monoidal category, natural in pointed functors $(X,x)\to(Y,y)$. Its nerve is a simplicial monoid, its component monoid is the group $\pi_1(BX,x)$, and its classifying space has the weak homotopy type of $\Omega_xBX$. Its multiplication agrees with the usual loop multiplication in $\Hot$. The construction uses no pushout-product axiom and no fibrant replacement of $X$.
\end{result}
\begin{proof}
Use concatenation and the empty path, as in C.1. A pointed functor applies to every vertex, arrow and natural-transformation component; it preserves concatenation and the empty path as equalities. The component identification follows from the path-word presentation of the groupoid completion. The space and its multiplication are identified by E.1, whose cited companion block treats composable pairs over $X^3$, so it identifies this operation and not just an abstract homotopy type. All constructions here use finite products, concatenation, interval collapses and filtered unions. No categorical pushout along a product of endpoint inclusions occurs.

Strict associativity supplies an $A_\infty$ multiplication on this model. The strict composition-compatible weak equivalence $\rho$ of B.2 identifies it as an $A_\infty$ loop space: here all arrows are weak, so the weak-vertical and full categories coincide, and the localization is the space $BX$. This higher identification uses B.2, while E.1 by itself supplies compatibility only in $\Hot$.
\end{proof}

For clarity, we fix the lifting-category convention. For $f:E\to S$, $e\in E$, and $\alpha:f(e)\to s$, an outgoing lift is an arrow $h:e\to e'$ with $f(e')=s$ and $f(h)=\alpha$; an arrow from $h$ to $h'$ is $k:e'\to e''$ with $f(k)=1_s$ and $kh=h'$. For $\alpha:s\to f(e)$, an incoming lift is $h:e'\to e$ over $\alpha$, with arrows $k:e'\to e''$ over $1_s$ satisfying $h'k=h$. In the $W_\infty$ sense, properness means that every outgoing lift category is weakly contractible, and smoothness means the analogous incoming condition. These are the lifting-category criteria of \citet[Definition 1.1 and Theorem 1.2]{Maltsiniotis2005Asphericite}, with duality fixing the variance. A Cartesian lift would in particular be terminal in its incoming lift category; a coCartesian lift would in particular be initial in its outgoing lift category.

\begin{result}{E.3}{Refuted}{The all-type endpoint functor is not a strict bifibration}
For $X=[1]=(0\xrightarrow w1)$, the functor $p:\Ch(X)\to X\times X$ is smooth in the $W_\infty$ lifting-category sense, but it is not proper and is neither a Grothendieck fibration nor a Grothendieck opfibration. In particular the duality argument asserting properness of this total category is false.
\end{result}
\begin{proof}
Smoothness is the incoming-lift contraction checked in the companion, VII ms.~164: an incoming lift is compared with a slice over the path obtained by adjoining the two given endpoint arrows; that slice has a terminal object. This works with arbitrary type maps in their stated direction.

For failure of properness take the empty path at $0$ and the outgoing base arrow $(1_0,w):(0,0)\to(0,1)$. An outgoing path morphism lies above a type map $e\to\tau$. Endpoint preservation forces $\tau=e$. Such a target cannot have endpoints $(0,1)$. The outgoing lift category is empty, which also excludes an opfibration.

For failure of a Grothendieck fibration take the empty path at $1$ and the incoming base arrow $(w,1_1):(0,1)\to(1,1)$. Its lift category is exactly $\Ch([1];0,1)$: a type collapse to $e$ is unique and the transformation to the constant path $1$ is unique. A Cartesian lift would be a terminal object $c$ in this category. Write $\tau$ for its type. Since its endpoints differ, $\tau$ is nonempty. Let $d$ be the path of type $\tau\tau$ which is constant at $0$ on the first block and is $c$ on the second. There are two different type projections $\tau\tau\to\tau$, collapsing respectively the first and the second block. For the first projection the pullback of $c$ is exactly $d$. For the second it is $c$ on the first block and constant at $1$ on the second. Pointwise inequalities give a natural transformation from $d$ to this second pullback; at the endpoints they are identities. Thus $d\to c$ has two distinct morphisms, contradicting terminality.

The program checks all $543$ such $[1]$-paths of length at most six and constructs these two morphisms; the argument proves the claim at every length. The finite base category has exactly two objects, its two identities and $w$, so all its composition data have been specified. Taking the opposite of a Grothendieck construction reverses its indexing category. Reversing a path inside a fixed $\Ty$ does not do this. That is why the claimed deduction of properness cannot be repaired by merely reversing $X$.
\end{proof}

\begin{result}{E.4}{Refuted}{Both homotopical lifting properties do not give Cartesian lifts}
The eventual alternating endpoint functor $\mathcal P(X)\to X\times X$ is smooth and proper for every $X$. Even for $X=[1]$ it is neither a Grothendieck fibration nor a Grothendieck opfibration.
\end{result}
\begin{proof}
The two homotopical lifting properties are precisely the even-even and odd-odd finite-support arguments checked in E.1. The relevant lift categories are filtered unions of categories with a terminal, respectively initial, object. Filtered unions preserve weak contractibility; they need not preserve those universal objects.

For $[1]$, consider incoming lifts into the constant path $1$ over $(w,1)$. Such a lift is an eventual path from $0$ to $1$, with its unique transformation to $1$. A terminal lift $d$ cannot exist: choose a single-transition path which changes from $0$ to $1$ across a direct edge so far to the left that $d$ is still $0$ just after that edge. Pointwise order then excludes a transformation from this path to $d$. Dually, outgoing lifts from the constant path $0$ over $(1,w)$ are the same eventual paths. For any proposed initial lift, choose a single-transition path which changes to $1$ later than that lift's eventual-$1$ tail begins. The required arrow to the later-transition path is impossible. The proof uses the two-object finite category just checked in E.3 and paths at arbitrarily large supports; a finite cut-off alone would miss the obstruction.
\end{proof}

\begin{result}{E.5}{Proved}{A functorial categorical mapping-path factorization}
For every functor $f:A\to X$, define
\[
 E(f)=A\times_{X,\sigma}\mathcal P(X),\qquad
 i_f(a)=(a,\operatorname{const}_{f(a)}),\qquad
 p_f(a,c)=\beta(c).
\]
Then $f=p_fi_f$, the functor $i_f:A\to E(f)$ is in $W_\infty$, and $p_f:E(f)\to X$ is smooth and proper. Its exact fibres compute the homotopy fibres of $Bf$. This factorization is natural in strictly commuting squares of functors. It is not asserted to be a Thomason cofibration--fibration factorization.
\end{result}
\begin{proof}
Smoothness and properness are stable under base change and composition: their incoming or outgoing lift categories for a pullback are the original lift categories, and the composite criterion follows by taking the appropriate comma construction. Thus the map $E(f)\to A\times X$, a base change of $p$, is smooth and proper, as is its composite with the projection $A\times X\to X$.

We justify the weak-equivalence step before using the section. A functor with both lifting properties has fibre-to-comma weak equivalences and satisfies Quillen's Theorem B, as in the checked proof CATHOT E.21. The same proof applies after every base change. It follows that any base-change square is homotopy Cartesian: compare its strict fibres with the derived fibres, and then apply the fibrewise weak-equivalence criterion for spaces over the same base. This includes empty fibres and empty bases. For $\sigma:\mathcal P(X)\to X$ all strict fibres are weakly contractible by E.1. Every base change of $\sigma$ is therefore a weak equivalence. In particular the projection $r:E(f)\to A$ is one. Since $ri_f=1_A$, two-out-of-three makes $i_f$ a weak equivalence. The homotopy-fibre description for $p_f$, and the equation $p_fi_f=f$, give the assertion about $Bf$.

For a commuting square $(u,v):f\to f'$, apply $u$ to $a$ and apply $v$ pointwise to the path; this defines the functor $E(f)\to E(f')$ and commutes strictly with both displayed structure maps. The usual categorical cocylinder factorization has the same purpose; \citet[Theorem 2.8]{Shoikhet2016} is an exact reference for Evrard's construction with its own indexing convention. We have not identified those conventions by an unproved isomorphism. The present proof needs no right properness of another chosen model structure and no product-of-cofibrations axiom.
\end{proof}

\begin{result}{E.6}{Refuted}{The standard eventual path object does not solve the five-element bifibration test}
Let $R=\{1,e_{00},e_{01},e_{10},e_{11}\}$ with
\[
 e_{ij}e_{kl}=e_{il},\qquad 1r=r1=r,
\]
and let $X=\mathbf B R$. Then $\mathcal P(X)\to X\times X$ is neither a Grothendieck fibration nor a Grothendieck opfibration. The all-type functor is not an opfibration either. These failures do not refute existence of some other strict bifibration representing the path local system over $X$.
\end{result}
\begin{proof}
Set $e=e_{00}$. Every product $er$ and $re$ is nonunit. Let $c$ be the constant diagram, whose edges are all $1$. An incoming lift of $(1,e)$ into $c$ consists of a path $d$ and a transformation $\alpha:d\to c$ with eventual components $1$ on the left and $e$ on the right. For an even vertex $i$ and an adjacent odd vertex $j$, its naturality equation is
\[
 \alpha_i=\alpha_jd_{i\to j}.
\]
For any sufficiently large odd $m$, there is such a lift with components $1$ for $i\leq m$ and $e$ for $i\geq m+1$; its sole nonidentity edge is $e$ on the arrow $m+1\to m$. Given any proposed terminal lift $\alpha$, choose a coordinate where $\alpha_i=e$, and choose $m$ beyond it. A morphism from this delayed lift $\gamma$ to $\alpha$ would require $\alpha_i\theta_i=\gamma_i=1$, which is impossible. Thus no terminal incoming lift exists.

For outgoing lifts $\alpha:c\to d$ the equation is $d_{i\to j}\alpha_i=\alpha_j$. A delayed step at an even $m$, with edge $e$ on $m\to m+1$, supplies another lift. If $\alpha$ were initial, the arrow to a sufficiently delayed lift would require $\theta_i\alpha_i=1$ at a coordinate where $\alpha_i=e$, again impossible. Finally an outgoing morphism from the empty path in the all-type category must have empty target; its two endpoint components coincide. Hence the base arrow $(1,e)$ has no such lift.

The program checks all $125$ associativity equations in $R$, all its nonunit-ideal equations, and the two families of naturality equations. Fiedorowicz's unnumbered lemma on pp.~33--34 identifies $BX\simeq S^2$ \citep{Fiedorowicz2002}; we opened and checked its five-element presentation. Our failure of strict lifts uses only the displayed multiplication table. It therefore addresses the precise standard path models in the STRATA E.11 test case, while leaving that existence question open.
\end{proof}

\subsection{The limit of the interval argument}

\begin{result}{E.7}{Refuted}{The pushout-product property fails already for the shortest alternating block}
For every $n\geq1$, let $\Phi_n=(0\to1\leftarrow2\to\cdots\leftarrow2n)$. The pushout-product
\[
 \alpha_n:\bigl(\{0\}\times\Phi_n\bigr)
       \amalg_{\{0\}\times\{0,2n\}}
       \bigl([1]\times\{0,2n\}\bigr)
       \longrightarrow [1]\times\Phi_n
\]
is a weak equivalence of finite categories but is not an intrinsic $W_\infty$-cofibration. For $n=1$ there is an explicit weak equivalence in its undercategory whose cobase change has source contractible and target of homotopy type $S^3$.
\end{result}
\begin{proof}
Put $P_n=[1]\times\Phi_n$ and write $l_i=(0,i)$, $u_i=(1,i)$. The source $A_n$ of $\alpha_n$ is the full sieve containing all $l_i$ and only $u_0,u_{2n}$. Its nerve is a tree, whereas $BP_n$ is contractible. Thus $\alpha_n\in W_\infty$.

We specify the pushout rather than treating it as a homotopy pushout. Let $C$ have objects $x,z,a,b$, arrows $\lambda,\mu:x\rightrightarrows z$, $r:z\to a$, $v:z\to b$, and the equal composites
\[
 r\lambda=r\mu=\eta:x\to a,\qquad
 v\lambda=v\mu=\theta:x\to b.
\]
There are no other nonidentity arrows. Its full subcategories on $x,z,a$ and on $x,z,b$ are cones on the parallel-arrow circle. Their nerves cover $NC$ and intersect in that circle, so $BC\cong S^2$. Let $D=C^{\triangleright}$, with new terminal object $s$. Define $F:A_n\to C$ by
\[
 F(l_0)=z,\quad F(u_0)=b,\quad F(w)=a\text{ at every other object}.
\]
The arrows $l_0\to l_1$ and $l_0\to u_0$ go to $r,v$, and all other nonidentity arrows of $A_n$ go to identities. The composite $A_n\to D$ is a weak equivalence since both categories are contractible.

For a sieve pushout $Q=P_n\amalg_{A_n}D$, the objects are those of $D$ and the remaining $u_k$, $1\leq k\leq2n-1$. Its crossing arrows are the explicit coends
\[
 \Hom_Q(d,u_k)=\int^{w\in A_n}
       \Hom_D(d,Fw)\times\Hom_{P_n}(w,u_k).
\]
There are no arrows from an upper complementary object to $D$. For $k=1$ the predecessor diagram maps finally to the fork $z\to a,z\to b$. For $n=1$ it has one extra object $u_2$ attached by $l_2\to u_2$, which does not change that finality. The fork inclusion into $C$ is ordinary final: at $x$ its comma category is the connected four-edge square, and at the other objects it has an initial object. The crossing set is consequently a singleton for every $d\in C$, and empty for $d=s$. For $k\geq2$ the predecessor diagram is nonempty connected and constant at $a$, so the crossing presheaf is $D(-,a)$. The upper transition to $u_1$ gives the unique map from that representable presheaf to the stated singleton presheaf on $C$.

It follows that $Q$ retracts onto the full subcategory $C\cup\{s,u_1\}$ by sending $u_k$, $k\geq2$, to $a$. The arrows represented by $1_a$ give a natural transformation from the inclusion after this retraction to $1_Q$. The retract consists of two terminal cones on $C$, with no arrow between their tips. Therefore
\[
 BQ\simeq\Sigma BC\simeq S^3.
\]
Cobase change along $\alpha_n$ has taken the weak equivalence $A_n\to D$ to $P_n\to Q$, which is not a weak equivalence. This is exactly failure of the defining preservation condition for an intrinsic cofibration.

The program constructs the coend equivalence relations, all arrows and all compositions, and checks every associativity equation for $n=1,2,3$. For $n=1$, $P_1$ has $(6,15)$ objects/arrows and $Q$ has $(6,20)$; the normalized rational Betti numbers of $Q$ are $(1,0,0,1,0)$. The cone proof establishes its integral homotopy type; numerical ranks alone would not do so. The supplied POLES H.6 states the obstruction for $n\geq2$. The extra predecessor at $n=1$ has been handled above, so the shortest block is included here. This strengthens the checked range without using the missing POLES computation files.
\end{proof}

\begin{result}{E.8}{Proved}{Finite interval cotensors have the right total type but need not have the right strict fibres}
For every finite oriented interval $I$, the constant-diagram map $X\to\Fun(I,X)$ is a categorical homotopy equivalence, natural in $X$. This does not make its endpoint evaluation a homotopy path fibration. In particular, the strict fibre of $\Fun([1],[1])\to[1]\times[1]$ at $(1,0)$ is empty, whereas the topological path space between those vertices is contractible and nonempty.
\end{result}
\begin{proof}
Order the interval's vertices geometrically as $0,\ldots,r$. The functors $v_k(i)=\min(i,k)$ join the constant-at-$0$ functor to $1_I$ by natural transformations oriented according to the successive edges. Precomposition gives the desired finite zigzag on $\Fun(I,X)$, relative to evaluation at $0$. Naturality in $X$ is immediate. For $I=X=[1]$ the indicated strict fibre consists of arrows $1\to0$, of which there are none. On realization $B[1]$ is an interval, and its endpoint path space is nonempty contractible. Thus the total equivalence cannot be restricted to arbitrary strict fibres. The eventual alternating construction in E.1 supplies precisely the additional relative result. Cotensoring and concatenation require no pushout-product axiom; invoking a monoidal model axiom for them would add a hypothesis not used in their proofs.
\end{proof}

\subsection{Other localizers and the proposed stratified interpretation}

\begin{result}{E.9}{Refuted}{Weak-equivalence preservation under pullback depends on the localizer}
The $W_\infty$-smooth and proper path functors cannot in general be called universally weak-equivalence-preserving fibrations for an arbitrary larger fundamental localizer. This fails for rational homology equivalences, even over the one-object category $\mathbf B C_2$.
\end{result}
\begin{proof}
Put $X=\mathbf B C_2$, with its unique object $x$, and let $q:\mathcal P(X;x,?)\to X$ be endpoint evaluation. Its total category is weakly contractible by E.1, and its strict fibre is of type $\Omega BC_2$, which has two contractible components. The functor $e\to X$ induces an isomorphism on rational homology. Indeed positive-degree group homology of a finite group with rational coefficients vanishes by the averaging contraction, and $BC_2$ is connected. In the pullback square along $q$, the induced map from the strict fibre to the contractible total category is not a rational homology equivalence: it maps $\mathbb Q^2$ to $\mathbb Q$ on $H_0$. The base category and group law are completely specified by $C_2=\{1,g\}$ and $g^2=1$.

This is a failure of universal preservation of a chosen larger weak-equivalence class. It does not contradict the contractibility of the incoming and outgoing lift categories, or the space-valued homotopy-fibre theorem. Those are the senses of smoothness, properness and homotopy fibration used in E.1--E.5. Any coefficient-specific claim needs its own definition and hypothesis.
\end{proof}

\begin{result}{E.10}{Refuted}{Components of $\Sigma$ do not generally define the proposed exit-path stratification}
For an arbitrary relative category $(M,\Sigma)$, there need not be a stratification of $BM$ whose strata are the components of the groupoid completion of $\Sigma$ and whose exit-path mapping spaces are either of the proposed relative path models.
\end{result}
\begin{proof}
Take the one-object category with endomorphism monoid $\{1,e\}$, $e^2=e$, and $\Sigma=\{1\}$. This is a saturated two-out-of-three class: its localization is the same ordinary category, and $e$ is not invertible. There is only one proposed stratum. The classifying space is contractible, since the natural transformation with component $e$ joins the identity functor to the constant functor sending every arrow to $1$. A one-stratum exit-path category is its fundamental $\infty$-groupoid, whose mapping space here is contractible. The relative localization instead has the discrete two-element endomorphism space $\{1,e\}$. Part B's isomorphism-only case gives the same two components for both path models. Hence neither can be that one-stratum exit-path space.

The obstruction is structural: within a stratum all morphisms of an exit-path category are invertible, while a relative localization can retain noninvertible endomorphisms. Positive statements for posets or suitable EI categories need additional hypotheses and an actual stratified comparison. The finite directed test with two parallel arrows has two components when only isomorphisms are inverted; inverting both arrows gives the ordinary circle path space with components indexed by $\mathbb Z$. This is consistent with the endpoint cases, but it does not justify the general stratification proposal.
\end{proof}

\begin{result}{E.11}{Open}{Existence of a different strict bifibration in the rectangular-band case}
For $X=\mathbf B R$ of E.6, we do not settle whether the path local system is represented by some strict bifibration over $X$ with weakly contractible total category and fibres of type $\Omega BX$. E.6 rules out the standard eventual endpoint construction as such a bifibration, and rules out the all-type endpoint construction as an opfibration. It leaves the existence question STRATA E.11 intact.
\end{result}
\begin{proof}[The remaining step]
For bases consisting entirely of monomorphisms or entirely of epimorphisms, the supplied STRATA construction uses an injectively fibrant diagram and adjoints to transition functors on categories of bounded contractible objects. In the rectangular-band base the four nonunits have neither cancellation property. Our path constructions give contractible lifting categories, not universal lifts. Replacing those categories by equivalent ones cannot be assumed to produce both adjoints coherently over the monoid. The missing step is precisely such a functorial, two-sided adjoint transition system or an obstruction to its existence.
\end{proof}

\begin{result}{E.12}{Known}{Contractible origin-path categories need not be filtered or totally aspherical}
For $Y=(z\to a,z\to b)$ and origin $z$, the exact all-type category $P(Y,z)$ is contractible in $\Cat\,\mathrm{hot}$, but is neither filtered nor totally $W_\infty$-aspherical. More generally $P(Y,y_0)\to Y$ is ordinary final when $BY$ is connected and simply connected.
\end{result}
\begin{proof}[Verification of TOTASPH G.1--G.6]
We checked the supplied proof with its exact type maps. For an initial origin, replacing every path by the constant path of the same type gives natural transformations to the identity and to the empty-path constant functor; hence the asserted two-step contraction. The paths ending at $a$ and at $b$ have no common successor, because their endpoint images would need one in $Y$.

For finality, an object of $y\backslash P(Y,y_0)$ is moved into the exact fibre over $y$ by appending its structure arrow backwards and collapsing that last edge. If $BY$ is simply connected this fibre is nonempty connected by the groupoid-completion computation of components. For this particular $Y$, its full inclusion in the four-object category $C$ of E.7 is ordinary final, with the connected square as comma category at $x$. Composing the two final functors gives a final image of $P(Y,z)$ whose classifying space is $S^2$. This violates total asphericity. The finite category $C$ and its complete composition table are independently checked in the program for E.7. Thus XIIba-C005 is false, with no claim that contractibility alone is insufficiently verified.
\end{proof}

\begin{result}{E.13}{Refuted}{Finite set-valued tests are not faithful}
Choose a small skeleton $\mathcal A_{\mathrm{fin}}$ of the full subcategory of $\Hot$ spanned by finite categories (finitely many objects and arrows). The restricted set-valued Yoneda functor
\[
 R:\Hot\longrightarrow
 \Fun(\mathcal A_{\mathrm{fin}}^{\op},\mathrm{Set}),
 \qquad R(Y)(A)=\Hom_{\Hot}(A,Y)\cong[BA,BY],
\]
is not faithful. An actual endofunctor of a small one-object category induces a map distinct from the identity in $\Hot$ but the identity on every one of these test sets. Thus the faithfulness argument in VIIc-C004 is false, although the particular equality $T^2=1$ is established independently in E.1.
\end{result}
\begin{proof}
We independently verify the finite-support argument of CATHOT F.2 \citep{SuppliedAuditsPATHS}. Let $G$ be the group of finite-support permutations of $\mathbb N=\{0,1,\ldots\}$ and define $s:G\to G$ by
\[
 s(\sigma)(0)=0,\qquad s(\sigma)(i+1)=\sigma(i)+1.
\]
This homomorphism induces the actual endofunctor $\mathbf B s:\mathbf BG\to\mathbf BG$ of the one-object category $\mathbf BG$; write $Bs:BG\to BG$ for its realization. Its image fixes $0$, whereas the transposition $(0\ 1)$ does not. Thus $s$ is not surjective, so $Bs$ is not a homotopy equivalence and cannot equal $1_{BG}$ in the homotopy category.

Every component of the nerve of a finite category has finitely generated fundamental group, since its one-skeleton is finite. If $P$ is a finitely generated group and $\rho:P\to G$ is a homomorphism, the supports of the images of a finite generating set have finite union. Choose $m$ with this union contained in $\{0,\ldots,m\}$. Conjugation by the finite cycle $c=(0\ 1\ \cdots\ m+1)$ satisfies
\[
 c\rho(p)c^{-1}=s(\rho(p))\qquad(p\in P).
\]
For a connected CW complex $Z$, maps to the $1$-type $BG$ are classified by homomorphisms $\pi_1Z\to G$ modulo conjugation. Indeed a based map is determined by its homomorphism on the one-skeleton, extends over the two-cells precisely when the relations hold, and extends uniquely up to homotopy over higher skeleta because the higher homotopy groups of $BG$ vanish; forgetting the basepoint introduces conjugation. Consequently $Bs\circ f\simeq f$ for every $f:BA\to BG$ with $A$ connected and finite. Apply this separately on every component for arbitrary finite $A$; the empty test has one map. Hence $R(\mathbf B s)=R(1_{\mathbf BG})$, although the maps are distinct in $\Hot$. The counterexample category $\mathbf BG$ is infinite; the argument treats every finite-category test.
\end{proof}
